\documentclass[aspectratio=169,11pt]{beamer}
% docs/slides/preamble.tex -- 五份 deck 共享
\documentclass[aspectratio=169,11pt]{beamer}

% ==== 中文与字体 ====
\usepackage[UTF8,fontset=none]{ctex}
\usepackage{fontspec}

\IfFontExistsTF{PingFang SC}{%
  \setCJKmainfont{PingFang SC}[BoldFont=PingFang SC,ItalicFont=PingFang SC,AutoFakeSlant=0.15]
  \setCJKsansfont{PingFang SC}[ItalicFont=PingFang SC,AutoFakeSlant=0.15]
  \setCJKmonofont{PingFang SC}
}{%
  \IfFontExistsTF{Source Han Sans SC}{%
    \setCJKmainfont{Source Han Sans SC}
    \setCJKsansfont{Source Han Sans SC}
  }{%
    \IfFontExistsTF{Noto Sans CJK SC}{%
      \setCJKmainfont{Noto Sans CJK SC}
      \setCJKsansfont{Noto Sans CJK SC}
    }{%
      \setCJKmainfont{Heiti SC}
      \setCJKsansfont{Heiti SC}
    }
  }
}

\IfFontExistsTF{Helvetica Neue}{\setmainfont{Helvetica Neue}}{%
  \IfFontExistsTF{Inter}{\setmainfont{Inter}}{}}
\IfFontExistsTF{JetBrains Mono}{\setmonofont{JetBrains Mono}[Scale=0.85]}%
  {\setmonofont{Menlo}[Scale=0.85]}

% ==== 颜色 ====
\usepackage{xcolor}
\definecolor{cMain}{HTML}{1F4E79}
\definecolor{cAccent}{HTML}{E07A1F}
\definecolor{cGray}{HTML}{595959}
\definecolor{cMute}{HTML}{F4F4F4}
\definecolor{cOK}{HTML}{2E7D32}
\definecolor{cBad}{HTML}{B23A3A}
\definecolor{cWarn}{HTML}{C8A300}

% ==== 主题 ====
\usepackage{tcolorbox}
\usetheme{metropolis}
\metroset{progressbar=frametitle,numbering=fraction,block=fill}
\setbeamercolor{frametitle}{bg=cMain,fg=white}
\setbeamercolor{progress bar}{fg=cAccent,bg=cMute}
\setbeamercolor{alerted text}{fg=cAccent}
\setbeamercolor{title}{fg=cMain}

% metropolis 默认尝试 Fira Sans，未安装时 fallback 到 Latin Modern Sans。
% 这里在主题加载后显式覆盖，强制使用 macOS 上一定存在的 Helvetica Neue，
% 跟中文 PingFang 风格一致；找不到则按顺序尝试 Inter / Arial。
\IfFontExistsTF{Helvetica Neue}{%
  \setsansfont{Helvetica Neue}%
}{%
  \IfFontExistsTF{Inter}{\setsansfont{Inter}}{%
    \IfFontExistsTF{Arial}{\setsansfont{Arial}}{}}%
}

% ==== TikZ ====
\usepackage{tikz}
\usetikzlibrary{positioning,arrows.meta,fit,shapes.geometric,calc,%
  decorations.pathmorphing,backgrounds}

\tikzset{
  node-box/.style={draw=cMain,thick,rounded corners=2pt,
    minimum height=8mm,minimum width=18mm,align=center,font=\small},
  node-state/.style={draw=cMain,thick,circle,minimum size=14mm,
    align=center,font=\small},
  node-ok/.style={draw=cOK,thick,fill=cOK!10},
  node-warn/.style={draw=cWarn,thick,fill=cWarn!15},
  node-bad/.style={draw=cBad,thick,fill=cBad!10},
  % 色盲友好的状态机三态：蓝 (正常/Closed) / 橙 (告警/Open) / 灰 (中间/HalfOpen)
  % 避免 red-green 二元区分，改用 hue + lightness 双通道。
  node-state-normal/.style={draw=cMain,thick,fill=cMain!12},
  node-state-alert/.style={draw=cAccent,thick,fill=cAccent!18},
  node-state-transition/.style={draw=cGray,thick,fill=cGray!18},
  arrow/.style={-{Stealth[length=2mm]},thick,draw=cMain},
  arrow-async/.style={-{Stealth[length=2mm]},thick,dashed,draw=cGray},
  arrow-bad/.style={-{Stealth[length=2mm]},thick,draw=cBad},
  arrow-alert/.style={-{Stealth[length=2mm]},thick,draw=cAccent},
  lifeline/.style={dashed,thick,draw=cGray},
}

% ==== 代码 ====
\usepackage{listings}
\lstdefinelanguage{Go}{
  morekeywords={break,case,chan,const,continue,default,defer,else,
    fallthrough,for,func,go,goto,if,import,interface,map,package,range,
    return,select,struct,switch,type,var,nil,true,false,iota},
  morekeywords=[2]{string,int,int64,uint,uint64,bool,byte,error,context},
  sensitive=true,
  morecomment=[l]//,morecomment=[s]{/*}{*/},
  morestring=[b]",morestring=[b]`,
}
\lstdefinelanguage{LuaRedis}{
  morekeywords={and,break,do,else,elseif,end,false,for,function,if,in,
    local,nil,not,or,repeat,return,then,true,until,while},
  morekeywords=[2]{redis,call,KEYS,ARGV,tonumber,HGET,HSET,GET,SET,
    DECR,DECRBY,INCR,INCRBY,EXPIRE,PEXPIRE,ZADD,ZCARD,ZREMRANGEBYSCORE,
    EVAL,EVALSHA},
  sensitive=true,
  morecomment=[l]{--},
  morestring=[b]",morestring=[b]',
}

\lstdefinestyle{base}{
  basicstyle=\ttfamily\scriptsize,
  numbers=left,numberstyle=\tiny\color{cGray},
  keywordstyle=\color{cMain}\bfseries,
  keywordstyle=[2]\color{cAccent},
  stringstyle=\color{cOK},
  commentstyle=\color{cGray}\itshape,
  backgroundcolor=\color{cMute},
  frame=single,framerule=0pt,
  showstringspaces=false,
  breaklines=true,
  tabsize=2,
  columns=fullflexible,
}
\lstset{style=base}

% ==== 三线表与附录页码 ====
\usepackage{booktabs}
\usepackage{appendixnumberbeamer}

% ==== 页脚来源标注 ====
\newcommand{\srcnote}[1]{{\color{cGray}\scriptsize\itshape #1}}

% 关键论点框
\newcommand{\keypoint}[1]{%
  \begin{tcolorbox}[colback=cMute,colframe=cMain,boxrule=0.5pt,arc=2pt]%
  #1\end{tcolorbox}}


\title{流量治理：从大促 GMV 倒推的五层闸门}
\subtitle{业务视角下的 TokenBucket / SlidingWindow / CircuitBreaker / 异步削峰 / HTTP cache}
\author{gomall · 业务侧 deck (v2)}
\date{}

\begin{document}

\maketitle

\begin{frame}[shrink]{今日大纲}
\tableofcontents
\end{frame}

% ============================================================
\section{业务困局：流量治理到底治什么}
% ============================================================

\begin{frame}[shrink]{一个真实的双 11 场景：5 秒转圈背后是多少 GMV}
\begin{itemize}
  \item 0 点开抢，预估峰值 \textbf{100k QPS} 涌进 \texttt{/orders/create}。
  \item 同步链路顶不住：DB 连接池打满 $\to$ 用户看到\textbf{5 秒转圈}。
  \item 行业数据：每多 1 秒延迟，下单放弃率 $+$ 7\%；\textbf{5 秒 $\approx$ 38\% 放弃}。
  \item 算账：单平台 0 点 10 亿 GMV 预算 $\times$ 38\% = \textbf{3.8 亿损失}，比限流误伤几百用户严重几个数量级。
  \item 结论：流量治理不是"挡攻击"，是\textbf{用业务可接受的代价换 GMV 不雪崩}。
\end{itemize}
\end{frame}

\begin{frame}[shrink]{5 个角色，5 种"流量治理是什么"}
\begin{tabular}{@{}lll@{}}
\toprule
角色 & 体感 & 关心 \\
\midrule
C 端用户  & 看到"系统繁忙"        & 我能不能再点 / 多久能买到 \\
商家     & 秒杀页打不开          & 我的店今天损失多少订单 \\
运营     & 0 点峰值能不能扛      & GMV 影响 / 黑产薅多少 \\
客服     & 70001 / 70002 投诉    & 怎么回话 / 怎么转 SRE \\
SRE     & 上游 429 / 下游慢     & SLO 是否还在预算内 \\
\bottomrule
\end{tabular}
\vspace{2mm}
\small 同一个 \texttt{70001} 业务码，对用户是"重试一下"、对客服是话术、对运营是\textbf{今天又少了多少单}、对 SRE 是 SLO 烧多少预算。\\
这份 deck 按这 5 个视角讲，不止讲算法。
\end{frame}

\begin{frame}[shrink]{核心哲学：宁可下游慢、不能上游 429}
\begin{itemize}
  \item README 写明的 SLO 取舍：P0 核心交易（\texttt{/orders/create} / \texttt{/paydown}）\textbf{99.95\%}（年宕 4h22min）、\textbf{不挂用户限流}。
  \item 业务翻译：\textbf{宁可让用户等 5 秒，也不能让用户看到 429}。
  \begin{itemize}
    \item 等 5 秒：体验差，但订单还在；
    \item 看到 429：用户立刻关 app，订单丢，怪平台不怪自己。
  \end{itemize}
  \item P0 路径靠\textbf{异步削峰 + 缓存 + 库存预扣}扛峰值，限流只挂全局兜底（IP 100 RPS）。
  \item P2 营销路径反过来：宁可误杀真人、不能让黑产薅光（\texttt{3/s/user} 激进阈值）。
  \item 这两套取舍分别声明在各领域的 routes.go 里（router.go 只做组合根），是流量治理"按业务等级分挂载"的核心。
\end{itemize}
\end{frame}

\begin{frame}[shrink]{gomall 五层闸门的业务定位}
\begin{tikzpicture}[node distance=4mm and 8mm]
  \node[node-box] (cli) {Client / CDN};
  \node[node-box,node-warn,right=of cli] (cdn) {HTTP cache\\卸读流量};
  \node[node-box,node-warn,right=of cdn] (tb) {TokenBucket\\IP 100 RPS};
  \node[node-box,below=of tb] (route) {router};
  \node[node-box,node-warn,below=of route] (sw) {SlidingWindow\\3/s/user};
  \node[node-box,node-warn,left=of sw] (cb) {CircuitBreaker\\5/10s/3};
  \node[node-box,node-warn,left=of cb] (mq) {OrderEnqueue\\ticket+MQ};
  \draw[arrow] (cli) -- (cdn);
  \draw[arrow] (cdn) -- (tb);
  \draw[arrow] (tb) -- (route);
  \draw[arrow] (route) -- (sw);
  \draw[arrow] (route) -- (cb);
  \draw[arrow] (route) -- (mq);
\end{tikzpicture}
\vspace{2mm}
\small 五层不是流水线串联，是\textbf{按接口按业务分级挑选挂载}：每层挡一个具体的角色痛点。
\end{frame}

\begin{frame}[shrink]{五层闸门 vs 五个角色：哪层解决谁的痛}
\footnotesize
\begin{tabular}{@{}lllll@{}}
\toprule
层 & 拦谁 & C 端体感 & 商家 / 运营关心 & 客服 \\
\midrule
HTTP cache    & 重复读  & 详情秒开    & CDN 卸了 80\% 读 & --- \\
TokenBucket   & 爬虫脚本 & 几乎无感    & GMV 保护         & --- \\
SlidingWindow & 撞库脚本 & 慢一点没事  & 反黑产 / 抢券公平 & 70001 \\
CircuitBreaker & 慢依赖  & "1 分钟后再试" & 防雪崩 / 守 SLO   & 70002 \\
OrderEnqueue  & 写峰值  & 拿 ticket 等一会 & 双 11 不挂        & "已入队、稍候查询" \\
\bottomrule
\end{tabular}
\vspace{2mm}
\small 每层都对应一个"角色 + 业务码 + 话术"，不是孤立算法。
\end{frame}

\begin{frame}[shrink]{核心 P0 vs 营销 P2：阈值取舍完全不同}
\begin{tabular}{@{}llll@{}}
\toprule
等级 & 接口举例 & 限流姿态 & 熔断 \\
\midrule
P0 交易 & /orders/create, /paydown & 不挂用户限流 & 必须有 \\
P1 读   & /product/show, /carousels & HTTP cache & 不需要 \\
P2 秒杀 & /skill\_product/skill    & 用户 3/s 激进 & 不需要 \\
P2 红包 & /redpacket/claim         & 用户 3/s     & 不需要 \\
P3 信息 & /category/list, 轮播图   & 全局 + CDN   & 不需要 \\
\bottomrule
\end{tabular}
\vspace{2mm}
\small \textbf{P0 宁可下游慢、不能上游 429}；\textbf{P2 宁可误杀、不能让黑产薅光}。\\
两类接口用同一套阈值就两头不讨好：P0 损 GMV、P2 黑产薅干。
\end{frame}

% ============================================================
\section{第一层：TokenBucket — 商家和爬虫的边界}
% ============================================================

\begin{frame}[shrink]{业务困局：爬虫脚本一晚上抓走多少商品数据}
\begin{itemize}
  \item 商家最讨厌的事：\textbf{竞对爬虫}夜里 5 QPS 持续 8 小时，把全店 SKU、价格、库存全抓走。
  \item 不拦的代价：① 数据被对家拿去比价杀价；② 单 IP 5 QPS 累计一夜 14 万次查询，DB 多扛 1 个商品详情连接池。
  \item 真用户单 IP 一晚上访问 \texttt{/product/show} 通常 $<$ 50 次（浏览几十个商品）。
  \item TokenBucket \textbf{100 RPS / 200 burst} 的业务含义：
  \begin{itemize}
    \item 真用户首屏并发 30+ 张图、瞬时 burst 100 $\to$ 200 桶吸收，\textbf{0 误伤}；
    \item 爬虫稳态 5 QPS 不超限，但触发\textbf{风控 + 业务监控告警}，配合 IP 段封禁。
  \end{itemize}
  \item 这层是\textbf{兜底闸门}，业务侧精细限流挂掉时至少不会被压死。
\end{itemize}
\end{frame}

\begin{frame}[shrink]{为什么 100 RPS / 200 burst：基线 64k 反推}
\begin{itemize}
  \item gomall 单机基线 \textbf{64{,}254 RPS}（\texttt{/ping}，stressTest \S 1）。
  \item 真实用户单 IP $<$ 1/600 单机容量；100 RPS 留够 600$\times$ 安全余量。
  \item burst $= 2 \times$ rps 的业务依据：
  \begin{itemize}
    \item 浏览器一个详情页 30+ 静态 GET 是常态；
    \item 商品详情图懒加载、瀑布流 $\to$ 瞬时 burst $\approx$ 100；
    \item 200 桶让一次正常加载\textbf{绝不触发}（业务零误伤）。
  \end{itemize}
  \item 商家视角：阈值定太松爬虫拿数据、定太死客户投诉；100/200 是\textbf{用 64k 基线 / 真实页面行为}双向算出的折中。
\end{itemize}
\end{frame}

\begin{frame}[shrink]{TokenBucket 算法：桶 + 令牌生成 + 漏出}
\begin{tikzpicture}[node distance=8mm and 14mm]
  \node[node-box,node-warn] (gen) {令牌生成\\100/s 稳定};
  \node[node-box,right=of gen] (bucket) {桶\\容量 200};
  \node[node-box,right=of bucket,node-ok] (out) {请求漏出};
  \node[node-box,below=of bucket,node-bad] (drop) {满了丢弃\\status=70001};
  \draw[arrow] (gen) -- (bucket);
  \draw[arrow] (bucket) -- (out);
  \draw[arrow-bad] (bucket) -- (drop);
\end{tikzpicture}
\vspace{2mm}
\begin{itemize}
  \item 令牌按 rps 匀速生成；桶满则溢出。
  \item 请求拿到令牌才放行，拿不到立刻 70001（不阻塞、不排队）。
  \item burst 决定能否吸收\textbf{瞬时尖刺}（首屏并发 / 用户狂点）。
  \item 业务含义：\textbf{允许尖刺 / 拒绝高水位}—— 真用户偶尔狂点不痛、爬虫稳态 5 QPS 触不到这层，但配合 IP 黑名单可以拦。
\end{itemize}
\end{frame}

\begin{frame}[shrink]{业务困局：限流器自己把整机内存吃光}
\begin{itemize}
  \item 每个新源 IP 第一次进来，都会在 \texttt{limiters} map 里新建一个 \texttt{rate.Limiter} 条目。
  \item 原实现\textbf{只增不删}：IP 用完就不管了，条目永久留在 map 里。
  \item 业务侧后果（长跑算账）：
  \begin{itemize}
    \item 公网攻击 / 扫描 / CDN 回源每天带来\textbf{百万级不重复 IP}；
    \item 一周累计上千万条目 $\to$ map 占用涨到 GB 级 $\to$ \textbf{进程 OOM 被内核杀掉}；
    \item 限流中间件\textbf{本该护住整机}，结果自己成了拖垮整机的那一个。
  \end{itemize}
  \item 隐蔽点：压测 / 灰度短时间看不出来，\textbf{只有真大促长时间高基数 IP 才暴雷}，最难复现。
  \item 修法：给每个条目记 \texttt{lastSeen}，后台 janitor 每分钟扫一遍，\textbf{踢掉超过 10 分钟没访问的}，让 map 容量恒定有界。
\end{itemize}
\end{frame}

\begin{frame}[fragile,shrink]{TokenBucket 中间件实现 (逐行)}
\begin{lstlisting}[language=Go,firstnumber=35,basicstyle=\ttfamily\tiny,
  caption={}]
func TokenBucket(rps rate.Limit, burst int) gin.HandlerFunc {
    var mu sync.Mutex
    limiters := make(map[string]*tokenBucketEntry) // 1) 带 lastSeen
    get := func(ip string) *rate.Limiter {
        mu.Lock(); defer mu.Unlock()
        e, ok := limiters[ip]
        if !ok {
            e = &tokenBucketEntry{limiter: rate.NewLimiter(rps, burst)}
            limiters[ip] = e
        }
        e.lastSeen = time.Now()                    // 2) 每次访问刷新
        return e.limiter
    }
    go func() {                                    // 3) 后台 janitor
        ticker := time.NewTicker(tokenBucketCleanupInterval)
        for range ticker.C {
            now := time.Now(); mu.Lock()
            for ip, e := range limiters {
                if now.Sub(e.lastSeen) > tokenBucketIdleTTL {
                    delete(limiters, ip)           // 4) 回收空闲条目
                }
            }
            mu.Unlock()
        }
    }()
    return func(c *gin.Context) { /* Allow() 不过 -> 200+70001 */ }
}
\end{lstlisting}
\scriptsize (1) 条目从裸 \texttt{*rate.Limiter} 换成带 \texttt{lastSeen} 的 \texttt{tokenBucketEntry}；(2) 每次取用都刷新 \texttt{lastSeen}（活跃 IP 永不被误删）；(3) 起一个 ticker goroutine 每分钟清扫；(4) 超过 \texttt{tokenBucketIdleTTL}（10min）未访问的条目删除 $\to$ \textbf{map 容量随活跃 IP 数有界、不再无界增长 OOM}；(5) 清扫和取用共用同一把锁，无竞态；(6) 进程内 map 是单机代价，多实例要换 Redis token bucket。
\end{frame}

% ============================================================
\section{第二层：SlidingWindow — 抢券公平与黑产 ROI}
% ============================================================

\begin{frame}[shrink]{业务困局：脚本党 10 万次刷券 = 真人少几张}
\begin{itemize}
  \item 秒杀 / 抢券场景：\textbf{业务对手是脚本}，不是用户。
  \begin{itemize}
    \item 真人最快点击 300--500ms / 次 $\to$ 1s 最多 2--3 次；
    \item 脚本不限速可以 1s 打 1000 次；
    \item 撞库账号一夜 24h$\times$3600s$\times$3 $=$ 26 万次抢券机会。
  \end{itemize}
  \item 不拦的业务后果（算账）：
  \begin{itemize}
    \item 10 万张优惠券，脚本 100 个号 $\times$ 1 千次 = 100 万次抢 $\to$ 99\% 命中率 = \textbf{9.9 万张被脚本抢走}；
    \item 真人只剩 \textbf{1\%}（1000 张）$\to$ 用户体感"根本抢不到"$\to$ 大促口碑崩。
  \end{itemize}
  \item SlidingWindow \texttt{3/s/user} 把脚本打回客户端，\textbf{反黑产 ROI 立等可现}。
\end{itemize}
\end{frame}

\begin{frame}[shrink]{阈值 3/s/user：用户 vs 脚本的精确分界}
\begin{itemize}
  \item 真用户行为画像：
  \begin{itemize}
    \item 双击 + 网络重试，瞬时 2 次同秒是常态；
    \item 留 50\% 余量 = \textbf{3 次 / 秒上限}，绝不误伤。
  \end{itemize}
  \item 脚本行为画像：第 4 次开始 70001，再脚本写多猛都白搭。
  \item TokenBucket 在此场景失效的原因：
  \begin{itemize}
    \item NAT 出口（学校 / 公司 / 4G 基站）下 \textbf{所有用户共享一个 IP}；
    \item 阈值放宽给真人 $\to$ 脚本一起放；阈值收紧拦脚本 $\to$ NAT 下真人被误伤。
  \end{itemize}
  \item 必须\textbf{按用户 ID}限，必须\textbf{分布式精确}（多实例下 Redis 共享状态）。
\end{itemize}
\end{frame}

\begin{frame}[shrink]{压测实测：3/s/user 阈值 2.2\% 误差}
\textbf{场景}：30 VU $\times$ 15s 打 \texttt{/skill\_product/skill}，共用同一 token（复现同一用户脚本的请求形态）。

\vspace{2mm}
\begin{tabular}{@{}lrr@{}}
\toprule
指标       & 实测      & 业务期望 \\
\midrule
通过请求数 & \textbf{46}     & 15s $\times$ 3 = 45 \\
被限流数   & 781{,}624  & 黑产全拦 \\
单接口 RPS & 52{,}082   & 与 ping 同量级 \\
误差       & \textbf{2.2\%} & $<$ 5\% \\
\bottomrule
\end{tabular}
\vspace{2mm}

\small 业务翻译：30 个脚本号刷了 78 万次，\textbf{真用户名额一张没少}；服务端只做一次 Lua RTT（52k RPS 同 ping 量级），\textbf{脚本成本被推到客户端}—— 这就是反黑产的 ROI。
\end{frame}

\begin{frame}[shrink]{SlidingWindow 在 Redis ZSet 上的实现}
\begin{tikzpicture}[node distance=6mm and 14mm]
  \node[node-box] (in) {请求};
  \node[node-box,node-warn,right=of in] (z1) {ZREMRANGE\\BYSCORE\\0..now-W};
  \node[node-box,right=of z1] (z2) {ZCARD\\$<$ Limit?};
  \node[node-box,node-ok,right=of z2] (z3) {ZADD\\now$\to$rand};
  \node[node-box,node-bad,below=of z2] (no) {0, count\\status=70001};
  \draw[arrow] (in) -- (z1);
  \draw[arrow] (z1) -- (z2);
  \draw[arrow] (z2) -- node[above,font=\tiny]{yes} (z3);
  \draw[arrow-bad] (z2) -- node[right,font=\tiny]{no} (no);
\end{tikzpicture}
\vspace{2mm}
\begin{itemize}
  \item 每个用户一个 ZSet：\texttt{rl:seckill:user:42}；
  \item score = 请求时间戳（ms）；member = "时间戳 + 6 字节随机"避免同毫秒冲突；
  \item 滑动窗口 = 每次请求先\textbf{踢掉 window 之外的旧记录}，再 ZCARD 数当前窗口内的请求数；
  \item 业务正确性：没有"窗口边界双倍突刺"，反黑产精度高。
\end{itemize}
\end{frame}

\begin{frame}[fragile,shrink]{SlidingWindow Lua 脚本 (逐行)}
\begin{lstlisting}[language=LuaRedis,firstnumber=18,basicstyle=\ttfamily\tiny,
  caption={}]
local key, window_ms = KEYS[1], tonumber(ARGV[1])
local limit, now_ms, member = tonumber(ARGV[2]), tonumber(ARGV[3]), ARGV[4]
redis.call('ZREMRANGEBYSCORE', key, 0, now_ms - window_ms)
local count = redis.call('ZCARD', key)
if count >= limit then return {0, count} end
redis.call('ZADD', key, now_ms, member)
redis.call('PEXPIRE', key, window_ms)
return {1, count + 1}
\end{lstlisting}
\small 逐行：① 入参 key/窗口/阈值/now/请求 id；② \textbf{ZREMRANGEBYSCORE} 滑掉过期（业务正确性的关键）；③ ZCARD 数当前窗口；④ 超阈值返回 0；⑤ 通过则 ZADD 入；⑥ PEXPIRE 自动回收（不留垃圾 key）；⑦ \textbf{整脚本单原子}，无竞态，是抢券公平的工程保证。
\end{frame}

\begin{frame}[shrink]{红包接口：同一套机制再用一次}
\begin{itemize}
  \item \texttt{/redpacket/claim} 也是 \textbf{Limit:3 Window:1s ByUser:true}（internal/redpacket/routes.go:20--25）。
  \item 与秒杀不同的是：红包额外挂\textbf{幂等}（避免一次抢包重复扣余额）。
  \item 业务模式归纳：
  \begin{itemize}
    \item \textbf{SlidingWindow = P2 营销标配}（反黑产 / 抢券公平）；
    \item \textbf{幂等 = P0/P2 资金链路标配}（防重复扣款）；
    \item 两个中间件\textbf{可以任意组合}，业务侧一行声明搞定。
  \end{itemize}
  \item 运营视角：双 11 / 春节 / 618 上新活动，复用现有中间件 = \textbf{零新代码、当天上线}。
\end{itemize}
\end{frame}

% ============================================================
\section{第三层：CircuitBreaker — 支付宕机不让交易链路死}
% ============================================================

\begin{frame}[shrink]{业务困局：第三方支付 30s 超时 = 全站雪崩}
\begin{itemize}
  \item 真实事故剧本（行业公开过的案例）：
  \begin{itemize}
    \item 13:42 第三方支付网关网络异常，RT 从 200ms 飙到 30s；
    \item 13:43 \texttt{/paydown} 全员 30s 超时 $\to$ goroutine 堆积 $\to$ 内存爆 $\to$ \textbf{gomall 自身挂}；
    \item 13:45 商家秒杀页打不开、订单详情打不开 $\to$ 客服炸锅；
    \item 13:50 SRE 才看到告警 —— 但人已经救不回来了，全站雪崩 40min。
  \end{itemize}
  \item 不熔断的代价：\textbf{下游慢 = 自己跟着挂}（连锁雪崩比单点宕机损失大 10$\times$）。
  \item CircuitBreaker 的业务含义：\textbf{我代下游回答你不行，免得你等 30 秒还白等}。
  \item P0 SLO 99.95\% 的预算 4h22min/年，雪崩一次烧完。
\end{itemize}
\end{frame}

\begin{frame}[shrink]{熔断保护的不是自己，是下游 + 自己 + 用户}
\begin{itemize}
  \item gomall \texttt{/paydown} 的下游 = 第三方支付（Stripe / 微信 / 支付宝）。
  \item 失败模式：\textbf{慢}（不是 500）。50x 可以快速判错，慢调用最毒。
  \item CircuitBreaker 三态阈值 \textbf{5 失败 / 10s 等待 / 3 探测}（internal/payment/routes.go:15--19）的业务含义：
  \begin{itemize}
    \item \textbf{5} = 容忍偶发抖动不一抖就跳 $\to$ 不误伤健康下游；
    \item \textbf{10s} = 给下游一次完整告警 + 自愈周期 $\to$ 配合下游 SRE 处理；
    \item \textbf{3} = HalfOpen 探测窗口够小 $\to$ 避免恢复瞬间再次打死。
  \end{itemize}
  \item Open 期间订单怎么办：落 outbox 等回路（详见 deck 11 Outbox），\textbf{不丢单}、不写脏数据。
\end{itemize}
\end{frame}

\begin{frame}[shrink]{CircuitBreaker 三态机}
\begin{tikzpicture}[node distance=18mm and 30mm]
  \node[node-state,node-state-normal] (closed) {Closed};
  \node[node-state,node-state-alert,right=of closed] (open) {Open};
  \node[node-state,node-state-transition,below=of open] (half) {HalfOpen};
  \draw[arrow-alert] (closed) -- node[above,font=\tiny]{连续 5 次失败} (open);
  \draw[arrow] (open.south west) to[bend right=20] node[left=2pt,font=\tiny,pos=0.5]{10s 后试探} (half.north west);
  \draw[arrow] (half) to[bend right=35] node[below,font=\tiny,pos=0.5]{3 个都成功} (closed);
  \draw[arrow-alert] (half.north east) to[bend right=20] node[right=2pt,font=\tiny,pos=0.5]{任一失败} (open.south east);
\end{tikzpicture}
\vspace{1mm}
\small \textbf{Closed}：放行所有，连续 5 次失败 $\to$ Open。
\textbf{Open}：直接返回 70002，10s 后试探。
\textbf{HalfOpen}：放 3 个探测，\textbf{3 个全成功}才 $\to$ Closed；任一失败立即 $\to$ Open。\\
非对称信任：\textbf{进 Open 容易、回 Closed 难}—— 抖动期不靠偶发成功放全量，是资金链路的正确取舍。
\end{frame}

\begin{frame}[shrink]{业务困局：半开期一个偶发成功就放全量流量}
\begin{itemize}
  \item 熔断 Open 后 10s 进 HalfOpen，放 3 个探测请求去问下游"你好了没"。
  \item 原判据有 bug：\textbf{只要其中一个探测成功，就立刻回 Closed} 放全量流量。
  \item 为什么这是灾难（抖动期算账）：
  \begin{itemize}
    \item 下游正在网络抖动 / 半死不活，成功率可能只有 30\%；
    \item 3 个探测里\textbf{碰巧第一个成功} $\to$ 误判"下游已恢复" $\to$ 全量 100k QPS 瞬间压回去；
    \item 半死的下游被这波全量\textbf{再次打垮} $\to$ 又熔断 $\to$ 又探测碰运气 $\to$ 反复横跳，永远恢复不了。
  \end{itemize}
  \item 正确判据：HalfOpen 期间\textbf{放行的探测必须全部成功}（连续成功达到 \texttt{HalfOpenMaxReq}）才回 Closed；任一失败立即回 Open。
  \item 修法：新增 \texttt{halfOpenSuccess} 计数，\textbf{只数成功的探测}，与"已放行数 \texttt{halfOpenReq}"分开统计。
\end{itemize}
\end{frame}

\begin{frame}[fragile,shrink]{CircuitBreaker.allow 状态转移 (逐行)}
\begin{lstlisting}[language=Go,firstnumber=77,basicstyle=\ttfamily\tiny,
  caption={}]
func (cb *circuitBreaker) allow() error {
    switch circuitState(cb.state.Load()) {
    case stateClosed: return nil                          // 1) 关闭态零开销
    case stateOpen:
        if time.Now().UnixNano()-cb.openedAt.Load() < int64(cb.opt.OpenTimeout) {
            return errCircuitOpen                         // 2) Open 期间一律拒
        }
        cb.mu.Lock(); defer cb.mu.Unlock()
        if circuitState(cb.state.Load()) == stateOpen && /* 3) 锁内 double-check */
            time.Now().UnixNano()-cb.openedAt.Load() >= int64(cb.opt.OpenTimeout) {
            cb.state.Store(int32(stateHalfOpen))
            cb.halfOpenReq.Store(0); cb.halfOpenSuccess.Store(0) // 4) 双计数清零
        }
        if circuitState(cb.state.Load()) == stateHalfOpen {
            if cb.halfOpenReq.Add(1) > cb.opt.HalfOpenMaxReq {
                return errCircuitOpen                     // 5) 探测名额超限退回
            }
            return nil
        }
        return errCircuitOpen
    case stateHalfOpen:
        if cb.halfOpenReq.Add(1) > cb.opt.HalfOpenMaxReq { return errCircuitOpen }
        return nil
    }
    return nil
}
\end{lstlisting}
\scriptsize (1) Closed 仅 atomic load（业务无感）；(2) Open 内不查时间就拒（10ms 内回 70002，对比 30s 超时是 3000$\times$ 提升）；(3) Open$\to$HalfOpen 锁内 double-check；(4) 进 HalfOpen 同时把\textbf{已放行数与成功数双双清零}（成功数是本轮修复新增）；(5) HalfOpen 只放 \texttt{HalfOpenMaxReq} 个探测，超额一律拒。
\end{frame}

\begin{frame}[fragile,shrink]{report：必须全部探测成功才回 Closed}
\begin{lstlisting}[language=Go,firstnumber=110,basicstyle=\ttfamily\tiny,
  caption={}]
func (cb *circuitBreaker) report(failed bool) {
    switch circuitState(cb.state.Load()) {
    case stateClosed:
        if failed {
            if cb.failures.Add(1) >= cb.opt.FailureThreshold {
                cb.tripOpen()             // 连续失败到阈值
            }
        } else { cb.failures.Store(0) }   // 成功重置
    case stateHalfOpen:
        if failed {
            cb.tripOpen()                 // 任一失败 -> 立即回 Open
        } else if cb.halfOpenSuccess.Add(1) >= cb.opt.HalfOpenMaxReq {
            cb.mu.Lock()                  // 成功数累计达阈值 -> 回 Closed
            cb.state.Store(int32(stateClosed))
            cb.failures.Store(0)
            cb.mu.Unlock()
        }
    }
}
\end{lstlisting}
\scriptsize (1) Closed 累加失败计数；(2) 成功重置避免偶发抖动跳闸；(3) \textbf{本轮修复点}：HalfOpen 回 Closed 的判据从"已放行数 \texttt{halfOpenReq} 到阈值"改成"\textbf{成功数 \texttt{halfOpenSuccess} 到阈值}"，即放行的探测\textbf{全部成功}才闭合；(4) 任一探测失败立即 \texttt{tripOpen} 回 Open；(5) "进 Open 容易、回 Closed 难"对资金链路合适：抖动期\textbf{绝不靠一个偶发成功就放全量}，避免半死下游被再次打垮。
\end{frame}

\begin{frame}[shrink]{熔断 + 幂等的三角配对：Open 返回 200 + 70002}
\begin{tikzpicture}[node distance=8mm and 14mm]
  \node[node-box] (cli) {客户端};
  \node[node-box,node-warn,right=of cli] (cb) {CircuitBreaker};
  \node[node-box,node-ok,right=of cb] (id) {Idempotency};
  \node[node-box,below=of id] (pay) {Stripe / 微信};
  \draw[arrow] (cli) -- (cb);
  \draw[arrow] (cb) -- (id);
  \draw[arrow] (id) -- (pay);
  \draw[arrow-bad,bend left=30] (cb.south) to node[below,font=\tiny]{Open: 200+70002} (cli.south);
\end{tikzpicture}
\vspace{1mm}
\begin{itemize}
  \item Open 返回 \textbf{200 + status=70002}（不是 5xx）：app 端老网关吞不了 5xx，统一业务码客户端只用一套 switch。
  \item 客户端拿 70002 指数退避（1s/2s/4s）；重试必须带原 Idempotency-Key。
  \item paydown 同时挂 CB + Idempotency 不是巧合：缺一边 $\to$ 重复扣款；缺另一边 $\to$ 雪崩。
\end{itemize}
\end{frame}

% ============================================================
\section{第四层：异步下单 — 双 11 0 点不转圈}
% ============================================================

\begin{frame}[shrink]{双 11 0 点的业务账：5 秒转圈 = 3.8 亿损失}
\begin{itemize}
  \item 同步下单链路：HTTP $\to$ ReserveStock $\to$ INSERT order $\to$ COMMIT $\to$ 回客户端，p99 约 50--100ms。
  \item 100k QPS $\times$ 100ms = \textbf{1 万并发}进 MySQL，连接池打爆 $\to$ p99 飙到 5s+。
  \item 用户行为模型：
  \begin{itemize}
    \item 等 1s 离开 5\%；
    \item 等 3s 离开 22\%；
    \item 等 5s 离开 \textbf{38\%}；
    \item 等 10s 离开 60\%+。
  \end{itemize}
  \item 单平台 0 点 10 亿 GMV $\times$ 38\% = \textbf{3.8 亿损失}—— 这是异步削峰的业务 ROI。
  \item gomall 方案：\textbf{把 DB 落库从同步链路拆出来}，前端拿 ticket 轮询。
\end{itemize}
\end{frame}

\begin{frame}[shrink]{用户体验业务量化：从"5 秒转圈"到"$<$ 10ms 拿 ticket"}
\begin{tabular}{@{}lll@{}}
\toprule
环节 & 同步 & 异步 ticket \\
\midrule
点击下单 & --- & --- \\
HTTP 响应 & p99 5s+ 转圈 & \textbf{$<$ 10ms 拿 ticket} \\
DB 落库 & 同步等 & MQ 慢慢落 \\
前端展示 & 5s 后看到结果 & 0.5s $\times$ 6 次轮询拿结果 \\
6 次后还没好 & 用户已离开 & 切"长等待"文案 \\
失败兜底 & 用户重试 = 重复下单风险 & 同 ticket 重查不下重 \\
\bottomrule
\end{tabular}
\vspace{2mm}
\small "等"这件事在异步路径里\textbf{从前台变后台}：用户先拿到"已入队"的确定性反馈，再轮询，再切长等待文案 —— 体感从"卡死"变成"在排队"，放弃率掉到个位数。
\end{frame}

\begin{frame}[shrink]{同步下单 vs 异步 ticket 链路}
\begin{tikzpicture}[node distance=4mm and 10mm]
  \node[node-box] (c1) {Client};
  \node[node-box,node-warn,right=of c1] (h1) {HTTP};
  \node[node-box,right=of h1] (db1) {Reserve+INSERT};
  \node[font=\tiny,below=2mm of db1] (s1) {同步：p99 100ms};

  \node[node-box,below=12mm of c1] (c2) {Client};
  \node[node-box,node-ok,right=of c2] (h2) {HTTP};
  \node[node-box,right=of h2] (mq) {Reserve+Ticket+MQ};
  \node[node-box,node-warn,below=4mm of mq] (cons) {Consumer\\慢慢落库};
  \node[node-box,below=4mm of h2] (poll) {GET status\\轮询 ticket};
  \draw[arrow] (c1) -- (h1); \draw[arrow] (h1) -- (db1);
  \draw[arrow] (c2) -- (h2); \draw[arrow] (h2) -- (mq);
  \draw[arrow-async] (mq) -- (cons);
  \draw[arrow] (c2) to[bend right=30] (poll);
\end{tikzpicture}
\vspace{2mm}
\small 上：所有客户端被 DB 阻塞；下：客户端 $<$ 10ms 拿到 ticket，DB 压力被 MQ 缓冲成\textbf{恒定 1k QPS 平流}。
\end{frame}

\begin{frame}[fragile,shrink]{OrderEnqueue：核心 12 行}
\begin{lstlisting}[language=Go,firstnumber=118,basicstyle=\ttfamily\tiny,
  caption={}]
func (s *OrderSrv) OrderEnqueue(ctx, req) (interface{}, error) {
    u, _ := ctl.GetUserInfo(ctx)
    if err := cache.ReserveStock(ctx, req.ProductID, req.Num); err != nil {
        return nil, err                                  // 1) 库存预扣
    }
    ticket := fmt.Sprintf("%d", snowflake.GenSnowflakeID()) // 2) ticket
    if err := defaultTicketStore.Put(ctx, ticket,
        OrderTicketStatus{Status: "pending"}, time.Hour); err != nil {
        _ = cache.ReleaseReservation(...); return nil, err   // 3) 写 Redis
    }
    body, _ := json.Marshal(AsyncOrderTask{Ticket: ticket, ...})
    if err := defaultAsyncProducer.Publish(ctx, body); err != nil {
        _ = cache.ReleaseReservation(...); return nil, err   // 4) 投 MQ
    }
    return OrderEnqueueResp{Ticket: ticket, Status: "pending"}, nil
}
\end{lstlisting}
\scriptsize ① 库存预扣前置，不足直接拒（用户立即看到"已售罄"50001 而非排队 1min 后才知道）；② snowflake 生成 ticket 无 DB；③ Redis pending 给前端轮询；④ 任一步失败 \textbf{ReleaseReservation} 回滚库存（Saga）；⑤ 投递成功才算入队，零丢单。
\end{frame}

% ============================================================
\section{第五层：HTTP cache — 让读流量根本不到 gomall}
% ============================================================

\begin{frame}[shrink]{业务困局：80\% 读流量都在重复问同一个问题}
\begin{itemize}
  \item 商品详情页：用户 A 看了一遍、用户 B 看了一遍、爬虫 C 看了一遍，\textbf{后端三次 DB 查询}。
  \item P1 P3 读接口（\texttt{/product/show}, \texttt{/carousels}, \texttt{/category/list}）数据更新频率\textbf{远低于查询频率}（运营每天调几次，用户每秒查几千次）。
  \item 卸到浏览器 / CDN 的业务收益：
  \begin{itemize}
    \item 浏览器命中：返回 304 Not Modified，\textbf{0 字节 body、0 DB 查询}；
    \item CDN 命中：根本不到 gomall，业务侧扩容压力直接砍 80\%；
    \item 用户体感：详情页\textbf{秒开}，首屏 LCP 提升 0.5s+。
  \end{itemize}
  \item gomall 实现：\texttt{HTTPCache(maxAge)} 中间件统一发 ETag + Cache-Control。
\end{itemize}
\end{frame}

\begin{frame}[shrink]{HTTP cache 的三种命中}
\begin{tikzpicture}[node distance=6mm and 12mm]
  \node[node-box] (c) {Client};
  \node[node-box,node-ok,right=of c] (cdn) {CDN};
  \node[node-box,node-warn,right=of cdn] (gw) {gomall};
  \node[node-box,below=4mm of gw] (db) {DB};
  \draw[arrow,bend left=40] (c.north) to node[above,font=\tiny]{1.强缓存命中 (max-age)} (c.north);
  \draw[arrow] (c) -- node[above,font=\tiny]{2.协商缓存} (cdn);
  \draw[arrow-async] (cdn) -- node[above,font=\tiny]{3.回源 ETag} (gw);
  \draw[arrow] (gw) -- (db);
\end{tikzpicture}
\vspace{2mm}
\begin{itemize}
  \item \textbf{强缓存}（Cache-Control: max-age=60）：浏览器本地命中，0 RTT $\to$ 用户秒开；
  \item \textbf{协商缓存}（ETag + If-None-Match）：到 CDN，命中 304 无 body $\to$ 省带宽；
  \item \textbf{回源}：CDN 不命中，回 gomall；ETag 用 SHA-256 前 16 字节算。
\end{itemize}
\end{frame}

\begin{frame}[shrink]{各接口的 max-age：运营改价的可接受延迟}
\begin{tabular}{@{}p{2.6cm}p{1.6cm}p{6.6cm}@{}}
\toprule
接口 & max-age & 业务依据 \\
\midrule
/product/list   & 30s  & 价格 / 库存可能变化，运营忍 30s 上下架 \\
/product/show   & 60s  & 商品详情更稳定 \\
/category/list  & 300s & 分类几乎不变 \\
/carousels      & 300s & 运营手动调，时效性弱 \\
\bottomrule
\end{tabular}
\vspace{2mm}\\
\small \textbf{max-age 的业务取舍}：太长 $\to$ "改价 / 下架"延迟生效（运营投诉）；太短 $\to$ 卸不了流量（SRE 投诉）。
300s = 运营改完最多 5min 生效，可接受；30s = 库存接近灵敏阈值的折中（敏感链路保护用户买不到错的）。
\end{frame}

% ============================================================
\section{客服话术：业务码就是产品}
% ============================================================

\begin{frame}[shrink]{业务码 70001 / 70002：客服话术对照}
\begin{tabular}{@{}lll@{}}
\toprule
业务码 & 业务含义 & 客服话术 \\
\midrule
70001 & 限流命中    & "活动太火爆，请稍等 1 秒重试" \\
70002 & 熔断开启    & "支付暂忙，1 分钟后再试，订单已保留" \\
60002 & 幂等命中    & "您已下单成功，无需重复" \\
50001 & 库存不足    & "已售罄" \\
\bottomrule
\end{tabular}
\vspace{3mm}
\small 同样是"请稍后"，三种含义完全不同：\\
\textbf{70001} 客户端 1s 重试就能过 $\to$ 用户自己解决；\\
\textbf{70002} 必须等下游自愈 + 客服转 SRE 看熔断面板 $\to$ 客服 + 工程协作；\\
\textbf{60002} 是用户自己重复了 $\to$ 客服安抚 + 引导查订单 $\to$ 客服自解决。\\
区分话术比统一一句"系统繁忙"对用户更友好，对客服也更省事。
\end{frame}

\begin{frame}[shrink]{70001 客诉：用户该怎么做 / 客服怎么转}
\begin{itemize}
  \item 用户视角：\textbf{"我点了下单，弹出'活动太火爆'是不是抢不到了？"}
  \item 客服话术：
  \begin{itemize}
    \item 第一句："您稍等 1 秒再点一次就好，不影响您的名额。"
    \item 第二句："这个提示是保护其他用户公平抢券的机制，您是真人不是脚本所以一定能过。"
    \item 不能说："系统出错了" —— 用户会觉得平台烂。
  \end{itemize}
  \item 客服后台看什么：
  \begin{itemize}
    \item 治理面板 \texttt{ratelimit:hit:seckill}，1 分钟内全平台命中数；
    \item 命中 $>$ 10$\times$ 历史均值 $\to$ 真大促，告诉用户"今天人多"；
    \item 命中正常但用户被拦 $\to$ 大概率脚本号 / 自动化插件。
  \end{itemize}
  \item 客服\textbf{不需要转 SRE}：70001 是设计内行为。
\end{itemize}
\end{frame}

\begin{frame}[shrink]{70002 客诉：客服 + SRE 双线联动}
\begin{itemize}
  \item 用户视角：\textbf{"我付钱付不出去，是不是要扣我两次？"}
  \item 客服话术：
  \begin{itemize}
    \item "支付通道暂时繁忙，1 分钟后再试，您的订单已保留 30 分钟不会过期。"
    \item "如果已经扣款但订单未付成功，请提供订单号，我们 10 分钟内确认。"
  \end{itemize}
  \item 客服后台动作：
  \begin{itemize}
    \item 看治理面板 \texttt{circuitbreaker:state:paydown} = \texttt{open}？
    \item \texttt{open} $\to$ 转 SRE，SRE 看下游通道（Stripe / 微信 / 支付宝），是否切备用；
    \item \texttt{closed} 但用户报 70002 $\to$ 老缓存 / 客户端旧版本。
  \end{itemize}
  \item 关键：\textbf{70002 = 设计内的"已知不可用"} $\neq$ "系统挂了"。客服话术要让用户感到\textbf{平台主动控制}，不是\textbf{失控}。
\end{itemize}
\end{frame}

% ============================================================
\section{大促运营 SOP：T-7 / T-1 / T 日 / T+1}
% ============================================================

\begin{frame}[shrink]{大促运营 SOP：流量治理的 T-7 / T-1 / T 日 / T+1}
\footnotesize
\begin{tabular}{@{}p{1.4cm}p{2.6cm}p{2.8cm}p{3.0cm}@{}}
\toprule
阶段 & 流量治理 & 客服 & SRE \\
\midrule
T-7 & 预设限流阈值 & 话术上线培训 & 容量评估 \\
    & 异步队列预拉  & 70001/70002 演练 & 副本数加倍 \\
T-1 & 灰度小流量 & 话术下发到坐席 & 监控面板待命 \\
    & MQ 队列清空  & 高频问题预备 & 报警分级核对 \\
T 日 0 点 & 实时观测命中 & 第一波处理 & 限流误伤 $<$ 5\% \\
       & CB 状态实时盯 & 转 SRE 工单  & SLO 烧不烧预算 \\
T+1 & 黑产盘点 & 客诉根因分析 & 复盘 / 调阈值 \\
    & 阈值复盘 & 高频客诉归档 & error budget 报告 \\
\bottomrule
\end{tabular}
\vspace{2mm}\\
\small SOP 的关键不是"工程很努力"，是\textbf{每个角色 T 日当天知道做什么、不做什么}。
\end{frame}

\begin{frame}[shrink]{大促容量评估：怎么算出"需要多少机器"}
\begin{itemize}
  \item 输入：① 业务侧预估的峰值 QPS；② 单机基线（gomall 64k RPS）；③ 安全系数（一般 1.5--2$\times$）。
  \item 计算：副本数 = $\lceil$ 预估峰值 $\times$ 安全系数 / 单机基线 $\rceil$。
  \item 双 11 例：100k QPS / 64k RPS $\times$ 1.5 = 2.34 $\to$ \textbf{至少 3 台}。
  \item 但秒杀 / 红包链路 \textbf{不是} 64k RPS 的水平：实测 52k 是 \texttt{/skill\_product/skill}，下单写链路 50k；要按\textbf{最弱链路}（这里是异步落库的 consumer 消费速度）算容量。
  \item 容量评估的常见错误：① 只看 RPS 不看 p99；② 没考虑下游放大（一个 paydown 触发 3 个下游调用）；③ 忽略冷启动（K8s 拉镜像 30s 起步）。
\end{itemize}
\end{frame}

% ============================================================
\section{SLO 与业务边界}
% ============================================================

\begin{frame}[shrink]{SLO：流量治理的"业务承诺"}
\begin{tabular}{@{}lll@{}}
\toprule
SLI & 阈值 & 业务含义 \\
\midrule
P0 接口可用率 & $\ge$ 99.95\%       & 年宕 $\le$ 4h22min（下单 / 支付） \\
限流误伤率   & $\le$ 5\%           & 真用户被 70001 命中 $<$ 5\% \\
熔断恢复 SLA & $\le$ 10s 探测 + 30s 全恢复 & 下游恢复后不滞后 \\
SlidingWindow 精度 & 误差 $\le$ 5\%      & 实测 \textbf{2.2\%} \\
基线 RPS    & $\ge$ 50k           & 实测 \textbf{64{,}254} \\
\bottomrule
\end{tabular}
\vspace{2mm}
\small SLO 是\textbf{用业务语言写给运营 / 商家 / SRE 看的承诺}，不是工程师内部 KPI。\\
每条 SLI 都对应一个"违约 = 怎么补救"的 runbook。
\end{frame}

\begin{frame}[shrink]{error budget：用数字管"什么时候停发版"}
\begin{itemize}
  \item P0 SLO $=$ 99.95\% $\to$ 月度\textbf{宕机预算 $\approx$ 21min}。
  \item 烧 1/3（7min）减发版；烧 2/3（14min）冻结非紧急；全烧光 $\to$ 下月转"稳定性 only"。
  \item 预算用 SLI 推：限流命中率 / 熔断打开时间 / 5xx 比例累加成"不可用时间"。
\end{itemize}
\keypoint{流量治理不是\textbf{一个机制}，是\textbf{五个机制按业务等级组合}。\\
每个阈值都要能用业务语言解释清楚。}
\end{frame}

\begin{frame}[shrink]{业务边界：流量治理不做什么}
\begin{itemize}
  \item \textbf{不做自适应限流}（Sentinel BBR）：阈值要能用业务语言算清楚，黑盒不接受；
  \item \textbf{不做按地域限流}：现阶段单区域，未做 GeoIP 维度切分；
  \item \textbf{不做黑产指纹库}：设备指纹 / 行为模型由独立风控团队（外部 L4）做；
  \item \textbf{不做 WAF / DDoS 防护}：L1 网络层由云厂商 Anti-DDoS 兜底；
  \item \textbf{不做商家自助调阈值}：merchant role 还在路线图，目前阈值平台配；
  \item \textbf{不做 risk score}：撞库 / 羊毛识别由风控团队接 70003 业务码（扩展位）。
\end{itemize}
\vspace{1mm}
\small 诚实列边界 $>$ 假装全能。\textbf{gomall 自己负责 L2 (IP) + L3 (User) + 写路径削峰}，挡住绝大部分业务异常流量。
\end{frame}

% ============================================================
\section{多层叠加 / 级联熔断 / 全链路压测}
% ============================================================

\begin{frame}[shrink]{多层流控漏斗：DDoS $\to$ IP $\to$ 用户 $\to$ 业务}
\begin{tikzpicture}[node distance=3mm]
  \node[node-box,node-warn] (l1) {L1 网络：CDN / Anti-DDoS（云厂商）};
  \node[node-box,node-warn,below=of l1] (l2) {L2 网关：IP TokenBucket（gomall）};
  \node[node-box,node-warn,below=of l2] (l3) {L3 业务：用户 SlidingWindow（gomall）};
  \node[node-box,node-warn,below=of l3] (l4) {L4 风控：设备指纹 / 行为模型（外部）};
  \draw[arrow] (l1) -- (l2);
  \draw[arrow] (l2) -- (l3);
  \draw[arrow] (l3) -- (l4);
\end{tikzpicture}
\vspace{1mm}
\begin{itemize}
  \item gomall 自己负责 L2 + L3，挡住\textbf{绝大部分业务异常流量}。
  \item L1 / L4 通常由云厂商和独立风控团队提供，不自己造。
\end{itemize}
\end{frame}

\begin{frame}[shrink]{服务调服务调服务：三层熔断的级联防护}
\begin{tikzpicture}[node distance=4mm and 10mm]
  \node[node-box] (cli) {Client};
  \node[node-box,node-warn,right=of cli] (api) {gomall API\\CB-A};
  \node[node-box,node-warn,right=of api] (svc) {订单服务\\CB-B};
  \node[node-box,node-warn,right=of svc] (pay) {支付聚合\\CB-C};
  \node[node-box,node-bad,below=of pay] (stripe) {Stripe};
  \node[node-box,node-bad,below=of svc] (alipay) {支付宝};
  \node[node-box,node-bad,below=of api] (wx) {微信支付};
  \draw[arrow] (cli) -- (api);
  \draw[arrow] (api) -- (svc);
  \draw[arrow] (svc) -- (pay);
  \draw[arrow] (pay) -- (stripe);
  \draw[arrow] (pay) -- (alipay);
  \draw[arrow] (pay) -- (wx);
\end{tikzpicture}
\vspace{1mm}
\begin{itemize}
  \item CB-C 守第三方：单家挂掉只断单家，其余通道照走。
  \item CB-B 守"支付聚合服务"整体：三家全挂时止血。
  \item CB-A 守"订单服务"整体：极端场景退化只读，让用户至少能看到订单列表。
\end{itemize}
\end{frame}

\begin{frame}[shrink]{级联防护的两个隐藏陷阱}
\begin{itemize}
  \item \textbf{陷阱 1：阈值同步雪崩}。如果 CB-A 与 CB-B 都设 5/10s，下游恢复瞬间所有上游同时 HalfOpen $\to$ 探测流量乘 N 倍打回去 $\to$ 下游再次挂。\\
  对策：外层 CB 的 OpenTimeout \textbf{要长于}内层（如 30s vs 10s），错开探测窗口。
  \item \textbf{陷阱 2：兜底链路也会挂}。CB-C Open 时若兜底返回的 cache stale 也过期，CB-B 会把 cache 读失败计作业务失败 $\to$ 跟着 Open。\\
  对策：兜底路径\textbf{不能}计入 CB-B 的失败计数；report 时按 status code 维度分桶。
  \item gomall 当前只挂了 CB-A（paydown），CB-B / CB-C 是把单实例 CB 复制到下游服务上的扩展模式。
\end{itemize}
\end{frame}

\begin{frame}[shrink]{网关层选型：Kong / APISIX / Nginx Ingress / 自研}
\begin{tabular}{@{}lllll@{}}
\toprule
方案 & 数据面 & 控制面 & 限流颗粒 & 适合 \\
\midrule
Nginx Ingress & nginx + Lua & K8s CRD & IP / Header & K8s 集群入口 \\
Kong          & nginx + Lua & PostgreSQL & 服务 / 消费者 & 中大型微服务 \\
APISIX        & nginx + Lua & etcd & 服务 / 路由 / Consumer & 高动态变更 \\
自研 gateway   & Go / Rust & 自管 & 业务字段任意 & 业务定制极重 \\
gomall 现状    & gin in-proc & 进程内 & IP / 用户 & 单体直连 \\
\bottomrule
\end{tabular}
\vspace{2mm}
\small gomall 没接独立网关，限流 / 熔断都跑在 gin 中间件里。\textbf{进生产前一般会接 Nginx Ingress + APISIX}：进程内中间件守业务码与状态机正确性，网关守 IP / DDoS / TLS。
\end{frame}

\begin{frame}[shrink]{五种主流限流算法：各自的业务边界}
\noindent\begin{tabular}{@{}lll@{}}
\toprule
算法 & 适合场景 & 局限 \\
\midrule
计数器 fixed-window     & 简单兜底 / 日级配额 & 窗口边界"双倍突刺" \\
滑动窗口 sliding-log    & 精确反黑产         & 内存 / Redis 开销大 \\
滑动窗口 sliding-counter & 折中               & 仍有边界误差   \\
令牌桶 token-bucket     & 允许突发           & 不抗持续高水位 \\
漏桶 leaky-bucket       & 强制平滑           & 突发延迟    \\
自适应 Sentinel BBR     & 容量自动反推       & 调参复杂 / 黑盒 \\
\bottomrule
\end{tabular}\par
\vspace{2mm}
{\small gomall 选\textbf{滑动窗口 ZSet}（sliding-log 变种）跑秒杀，因为业务上要\textbf{精确反黑产}；选\textbf{令牌桶}跑全局，因为允许首屏突发。}
\end{frame}

\begin{frame}[shrink]{熔断打开后：兜底数据让"不可用"变"降级可用"}
\begin{tikzpicture}[node distance=4mm and 10mm,every node/.append style={font=\scriptsize}]
  \node[node-box] (req) {请求};
  \node[node-box,node-warn,right=of req] (cb) {CB Open?};
  \node[node-box,node-ok,right=of cb] (live) {下游实时};
  \node[node-box,node-warn,below=of live] (stale) {cache stale $<$ 5min};
  \node[node-box,node-warn,below=of stale] (def) {默认值 (零、空数组)};
  \node[node-box,node-bad,below=of def] (rej) {拒绝 70002};
  \draw[arrow] (req) -- (cb);
  \draw[arrow] (cb) -- node[above,font=\tiny]{no} (live);
  \draw[arrow-bad] (cb) -- node[right,font=\tiny]{yes} (stale);
  \draw[arrow-bad] (stale) -- node[right,font=\tiny]{无 cache} (def);
  \draw[arrow-bad] (def) -- node[right,font=\tiny]{资金链路} (rej);
\end{tikzpicture}
\vspace{1mm}
\begin{itemize}
  \item 读路径（商品 / 推荐 / 评价）：cache stale $>$ 默认值 $>$ 拒绝。用户能看到旧数据胜过白屏。
  \item 写路径（支付 / 下单 / 扣库存）：\textbf{宁可拒，不要默认值}，否则资金 / 库存对不上账。
  \item 兜底数据来自 HTTPCache 中间件留下的 ETag 体，TTL 内可用。
\end{itemize}
\end{frame}

\begin{frame}[shrink]{业务码全表：错误码 $\to$ 客户端行为 $\to$ 客服话术}
\footnotesize
\begin{tabular}{@{}lllll@{}}
\toprule
码 & 含义 & 客户端 & 监控位 & 客服话术 \\
\midrule
20000 & 成功         & 正常        & 业务 KPI    & --                        \\
30001 & token 错误   & 跳登录      & 不告警      & "请重新登录"              \\
30002 & token 过期   & 静默续期    & 不告警      & --                        \\
30005 & 权限不足     & 不重试      & 不告警      & "无权限"                  \\
50001 & 库存不足     & 不重试      & 业务面板    & "已售罄"                  \\
60002 & 幂等命中     & 读结果      & 不告警      & "您已成功，无需重复"      \\
70001 & 限流命中     & 1s 重试     & 治理面板    & "活动太火、请稍后"        \\
70002 & 熔断开启     & 指数退避    & SRE 告警    & "支付暂忙、1 分钟后再试"  \\
70003 & 黑产识别命中 & 不重试      & 风控面板    & "为了账户安全、请联系客服" \\
\bottomrule
\end{tabular}
\vspace{1mm}
\end{frame}

% ============================================================
\appendix
% ============================================================

\begin{frame}[shrink]{面试 Q\&A (上)}
\textbf{Q1}：为什么 P0 下单不挂用户限流？\\
\hspace{4mm}\small A：README 写明的"宁可下游慢、不能上游 429"。下单看到 429 用户立刻关 app、订单丢且怪平台；等 5 秒虽然差但订单还在。P0 靠异步削峰 + 缓存 + 库存预扣扛峰值，限流只挂 IP 兜底。

\vspace{2mm}
\textbf{Q2}：阈值 3/s/user 是怎么定的？\\
\hspace{4mm}\small A：① 真人最快点击 300--500ms / 次 $\to$ 1s 最多 2--3 次，留 50\% 余量定 3；② 第 4 次开始 70001 截断脚本；③ 必须按 userID 不是 IP（NAT 出口共享 IP），多实例下用 Redis 共享状态。代码 \texttt{internal/skill/routes.go:17--24}。

\vspace{2mm}
\textbf{Q3}：CircuitBreaker 5/10s/3 三个数字有什么讲究？\\
\hspace{4mm}\small A：5 = 容忍偶发抖动不一抖就跳；10s = 给下游一次完整告警 + 自愈周期；3 = HalfOpen 探测够小，避免再次打死下游。三个数字都对应一个业务事故场景，不是工程师拍的。
\end{frame}

\begin{frame}[shrink]{面试 Q\&A (中)}
\textbf{Q4}：异步下单会不会丢单？怎么向运营解释 SLO？\\
\hspace{4mm}\small A：不会。库存先 ReserveStock（Redis Lua 原子），任一失败都 ReleaseReservation 回滚；ticket TTL=1h 兜底；MQ 用 RabbitMQ 持久化 + 手动 ack。代码 \texttt{internal/order/async.go:118--171}。运营 SLO：99.95\% 可用 = 年宕 4h22min，异步路径 p95 入队 $<$ 10ms、轮询 $<$ 3s 拿结果。

\vspace{2mm}
\textbf{Q5}：为什么熔断返回 200 不返回 503？客服怎么用？\\
\hspace{4mm}\small A：app 端老网关不能优雅处理 5xx；统一业务码 70002 = 客户端只用维护一套 status switch、客服直接对照话术表。监控不要按 4xx/5xx 算错误率，必须看 status。代码 \texttt{middleware/circuitbreaker.go:58--66}；客服话术 \texttt{pkg/e/msg.go:52}。

\vspace{2mm}
\textbf{Q6}：黑产薅券 ROI 怎么算 SlidingWindow 的收益？\\
\hspace{4mm}\small A：10 万张券，100 个脚本号 $\times$ 1 千次 = 100 万次抢，不拦 $\to$ 9.9 万张被脚本拿走、真人 1\%。挂 SlidingWindow 3/s/user $\to$ 脚本第 4 次 70001、78 万次被拦、真人名额一张没少。压测实测误差 \textbf{2.2\%}（\texttt{stressTest/REPORT.md \S 5}）。
\end{frame}

\begin{frame}[shrink]{面试 Q\&A (下)}
\textbf{Q7}：多层熔断会不会"探测雪崩"？怎么办？\\
\hspace{4mm}\small A：会。外层 CB 的 OpenTimeout 要长于内层（30s vs 10s）错开 HalfOpen 探测窗口；同时兜底路径不应计入失败计数，否则 cache 过期会让外层跟着 Open。代码 \texttt{middleware/circuitbreaker.go:70 失败判定}。

\vspace{2mm}
\textbf{Q8}：全链路压测怎么避免污染真实数据？\\
\hspace{4mm}\small A：流量染色：网关层注入 \texttt{X-Test} header $\to$ 透传到 ctx $\to$ DAO 看 ctx 切 \texttt{order\_shadow} 表；MQ 用独立 routing-key 走影子队列；Redis 用 \texttt{shadow:} 前缀隔离。设计稿 \texttt{middleware/coloring.go}（gomall 当前未接，扩展位）。

\vspace{2mm}
\textbf{Q9}：限流算法为什么不选自适应（Sentinel BBR）？\\
\hspace{4mm}\small A：自适应是黑盒、冷启动 5--10min 才稳；秒杀 3/秒/用户 这种阈值用业务语言能算清楚，运营 / 客服都看得懂，没必要让运维去看 BBR 曲线判断。\textbf{业务上能解释}比\textbf{工程上更聪明}更重要。代码 \texttt{repository/cache/ratelimit.go:18--34}。
\end{frame}

\begin{frame}[shrink]{代码位置一览}
\footnotesize
\begin{description}
  \item[router 挂载点]
  \item[TokenBucket]
  \item[SlidingWindow]
  \item[CircuitBreaker]
  \item[HTTPCache]
  \item[OrderEnqueue]
  \item[业务码]
  \item[压测报告]
  \item[README SLO 与边界]
\end{description}
\keypoint{业务 70\% / 代码 30\%：阈值都能用业务语言解释；代码都能在 file:line 翻到。}
\end{frame}

\end{document}
